% GENERATED FILE. Edit canonical lessons, then run npm run generate:books.
% canonical-source-hash: fnv1a64:d4cc3e3c0203a9a4
% canonical-lessons: HI-C426-raqam, HI-C426-naqad, HI-C426-cek, HI-C426-bailens, HI-C426-etiem

\chapter{Sums and Cash}
\label{ch:hi-sums-and-cash}
\hlchaptermodality{426}

\begin{chapteropening}
\textbf{By the end of this chapter:} \emph{I can pay in cash or by cheque, deposit a sum and find a cash machine.}

Complete the last lesson of chapter 426: I can pay in cash or by cheque, deposit a sum and find a cash machine.
\end{chapteropening}

\section[\texorpdfstring{\hi{रक़म} (raqam)}{रक़म (raqam)}]{\hi{रक़म} (raqam) — a sum, an amount (of money)}

\label{lesson:HI-C426-raqam}



\begin{quote}
Before the new one: say the Hindi for a stink, then the Hindi for to seep.
\end{quote}

\subsection*{You'll want to know: \hi{रक़म}}
\textbf{\hi{रक़म}} — \emph{raqam} — \textquotedblleft{}a sum, an amount (of money)\textquotedblright{}.

\begin{quote}
\hi{यह} \hi{रक़म} \hi{बहुत} \hi{बड़ी} \hi{है।}
\end{quote}

\begin{quote}
\emph{yah raqam bahut baṛī hai.}
\end{quote}

\begin{quote}
This is a very large sum.
\end{quote}

Feminine.

\begin{cousinweb}
\textbf{\hi{रक़म}} (\emph{raqam}) is usually traced to Arabic \emph{raqam}, \textquotedblleft{}a written figure, a mark\textquotedblright{}, through Persian: the figure written on a cheque, and so the sum itself.
\end{cousinweb}

\begin{grammarlens}[title={What you've built}]
One more word for getting around.
\end{grammarlens}

\subsection*{Your turn}
\begin{itemize}
\raggedright
  \item \emph{Say it:} \emph{raqam}
  \item \emph{Say it:} \emph{raqam}, once more
  \item \emph{Say it:} \emph{yah raqam bahut baṛī hai.} — This is a very large sum.
  \item \emph{Recall:} say the Hindi for a stink, then the Hindi for to seep, then say \emph{raqam} again
\end{itemize}

\begin{tcolorbox}[breakable,colback=teal!4,colframe=teal!35!black,title={Before you move on}]
What does \hi{रक़म} mean? (\textquotedblleft{}A sum, an amount (of money)\textquotedblright{}.) Say it once more.
\end{tcolorbox}
\section[\texorpdfstring{\hi{नक़द} (naqad)}{नक़द (naqad)}]{\hi{नक़द} (naqad) — in cash; cash}

\label{lesson:HI-C426-naqad}



\begin{quote}
Before the new one: say the Hindi for to seep, then the Hindi for a sum.
\end{quote}

\subsection*{You'll want to know: \hi{नक़द}}
\textbf{\hi{नक़द}} — \emph{naqad} — \textquotedblleft{}in cash; cash\textquotedblright{}.

\begin{quote}
\hi{मैंने} \hi{दुकानदार} \hi{को} \hi{नक़द} \hi{पैसे} \hi{दिए।}
\end{quote}

\begin{quote}
\emph{maiṁne dukāndār ko naqad paise diye.}
\end{quote}

\begin{quote}
I paid the shopkeeper in cash.
\end{quote}

\begin{cousinweb}
\textbf{\hi{नक़द}} (\emph{naqad}) is usually traced to Arabic \emph{naqd}, \textquotedblleft{}ready money\textquotedblright{}, through Persian; Hindi breaks the final qd with a vowel.
\end{cousinweb}

\begin{grammarlens}[title={What you've built}]
One more word for getting around.
\end{grammarlens}

\subsection*{Your turn}
\begin{itemize}
\raggedright
  \item \emph{Say it:} \emph{naqad}
  \item \emph{Say it:} \emph{naqad}, once more
  \item \emph{Say it:} \emph{maiṁne dukāndār ko naqad paise diye.} — I paid the shopkeeper in cash.
  \item \emph{Recall:} say the Hindi for to seep, then the Hindi for a sum, then say \emph{naqad} again
\end{itemize}

\begin{tcolorbox}[breakable,colback=teal!4,colframe=teal!35!black,title={Before you move on}]
What does \hi{नक़द} mean? (\textquotedblleft{}In cash; cash\textquotedblright{}.) Say it once more.
\end{tcolorbox}
\section[\texorpdfstring{\hi{चेक} (cek)}{चेक (cek)}]{\hi{चेक} (cek) — a cheque}

\label{lesson:HI-C426-cek}



\begin{quote}
Before the new one: say the Hindi for a sum, then the Hindi for in cash.
\end{quote}

\subsection*{You'll want to know: \hi{चेक}}
\textbf{\hi{चेक}} — \emph{cek} — \textquotedblleft{}a cheque\textquotedblright{}.

\begin{quote}
\hi{उसने} \hi{किराये} \hi{के} \hi{लिए} \hi{चेक} \hi{दिया।}
\end{quote}

\begin{quote}
\emph{usne kirāye ke liye cek diyā.}
\end{quote}

\begin{quote}
He gave a cheque for the rent.
\end{quote}

Masculine.

\begin{cousinweb}
\textbf{\hi{चेक}} (\emph{cek}) is borrowed from English \emph{cheque}.
\end{cousinweb}

\begin{grammarlens}[title={What you've built}]
One more word for getting around.
\end{grammarlens}

\subsection*{Your turn}
\begin{itemize}
\raggedright
  \item \emph{Say it:} \emph{cek}
  \item \emph{Say it:} \emph{cek}, once more
  \item \emph{Say it:} \emph{usne kirāye ke liye cek diyā.} — He gave a cheque for the rent.
  \item \emph{Recall:} say the Hindi for a sum, then the Hindi for in cash, then say \emph{cek} again
\end{itemize}

\begin{tcolorbox}[breakable,colback=teal!4,colframe=teal!35!black,title={Before you move on}]
What does \hi{चेक} mean? (\textquotedblleft{}A cheque\textquotedblright{}.) Say it once more.
\end{tcolorbox}
\section[\texorpdfstring{\hi{बैलेंस} (bailens)}{बैलेंस (bailens)}]{\hi{बैलेंस} (bailens) — a balance (in an account)}

\label{lesson:HI-C426-bailens}



\begin{quote}
Before the new one: say the Hindi for in cash, then the Hindi for a cheque.
\end{quote}

\subsection*{You'll want to know: \hi{बैलेंस}}
\textbf{\hi{बैलेंस}} — \emph{bailens} — \textquotedblleft{}a balance (in an account)\textquotedblright{}.

\begin{quote}
\hi{बैंक} \hi{में} \hi{मेरा} \hi{बैलेंस} \hi{कम} \hi{है।}
\end{quote}

\begin{quote}
\emph{baiṅk meṁ merā bailens kam hai.}
\end{quote}

\begin{quote}
My bank balance is low.
\end{quote}

Masculine.

\begin{cousinweb}
\textbf{\hi{बैलेंस}} (\emph{bailens}) is borrowed from English \emph{balance}, the money left in an account; \hi{बैलेंस} \hi{देखना}, to check your balance.
\end{cousinweb}

\begin{grammarlens}[title={What you've built}]
One more word for getting around.
\end{grammarlens}

\subsection*{Your turn}
\begin{itemize}
\raggedright
  \item \emph{Say it:} \emph{bailens}
  \item \emph{Say it:} \emph{bailens}, once more
  \item \emph{Say it:} \emph{baiṅk meṁ merā bailens kam hai.} — My bank balance is low.
  \item \emph{Recall:} say the Hindi for in cash, then the Hindi for a cheque, then say \emph{bailens} again
\end{itemize}

\begin{tcolorbox}[breakable,colback=teal!4,colframe=teal!35!black,title={Before you move on}]
What does \hi{बैलेंस} mean? (\textquotedblleft{}A balance (in an account)\textquotedblright{}.) Say it once more.
\end{tcolorbox}
\section[\texorpdfstring{\hi{एटीएम} (eṭīem)}{एटीएम (eṭīem)}]{\hi{एटीएम} (eṭīem) — a cash machine, an ATM}

\label{lesson:HI-C426-etiem}



\begin{quote}
Before the new one: say the Hindi for a cheque, then the Hindi for a balance.
\end{quote}

\subsection*{You'll want to know: \hi{एटीएम}}
\textbf{\hi{एटीएम}} — \emph{eṭīem} — \textquotedblleft{}a cash machine, an ATM\textquotedblright{}.

\begin{quote}
\hi{बाज़ार} \hi{के} \hi{पास} \hi{एक} \hi{एटीएम} \hi{है।}
\end{quote}

\begin{quote}
\emph{bāzār ke pās ek eṭīem hai.}
\end{quote}

\begin{quote}
There is a cash machine near the market.
\end{quote}

Masculine.

\begin{cousinweb}
\textbf{\hi{एटीएम}} (\emph{eṭīem}) is borrowed from English \emph{ATM}, said letter by letter.
\end{cousinweb}

\begin{grammarlens}[title={What you've built}]
One more word for getting around.
\end{grammarlens}

\subsection*{Your turn}
\begin{itemize}
\raggedright
  \item \emph{Say it:} \emph{eṭīem}
  \item \emph{Say it:} \emph{eṭīem}, once more
  \item \emph{Say it:} \emph{bāzār ke pās ek eṭīem hai.} — There is a cash machine near the market.
  \item \emph{Recall:} say the Hindi for a cheque, then the Hindi for a balance, then say \emph{eṭīem} again
\end{itemize}

\begin{tcolorbox}[breakable,colback=teal!4,colframe=teal!35!black,title={Before you move on}]
What does \hi{एटीएम} mean? (\textquotedblleft{}A cash machine, an ATM\textquotedblright{}.) Say it once more.
\end{tcolorbox}
